\section*{Introduction}
\addcontentsline{toc}{section}{Introduction}

% TODO: Rédiger l'introduction du chapitre 4
% Points à couvrir:
% 1. Objectif du chapitre (Réalisation/Implémentation)
% 2. Technologies utilisées
% 3. Plan du chapitre

[CONTENU À REMPLIR: Introduction du chapitre 4 - Réalisation]

\section{Environnement de développement}

% TODO: Décrire l'environnement de développement
% Points à inclure:
% - Langages de programmation
% - IDEs et outils
% - Serveurs
% - Gestion de version
% - Outils de test

\subsection{Technologies et outils utilisés}

[CONTENU À REMPLIR: Description des technologies et outils]

% Exemple:
% \begin{itemize}
%     \item \textbf{Frontend:} [Technologie et versions]
%     \item \textbf{Backend:} [Technologie et versions]
%     \item \textbf{Base de données:} [SGBD]
%     \item \textbf{IDE:} [IDE utilisé]
%     \item \textbf{Gestion de version:} [Système utilisé]
%     \item \textbf{Serveur:} [Type de serveur]
% \end{itemize}

\subsection{Configuration et installation}

[CONTENU À REMPLIR: Instructions de configuration et installation]

% Exemple:
% \begin{verbatim}
% 1. Cloner le repository
% 2. Installer les dépendances
% 3. Configurer les variables d'environnement
% 4. Initialiser la base de données
% 5. Lancer l'application
% \end{verbatim}

\section{Implémentation du frontend}

% TODO: Décrire l'implémentation du frontend
% Points à inclure:
% - Structure du projet
% - Composants développés
% - Fonctionnalités implémentées
% - Captures d'écran

\subsection{Structure du projet frontend}

[CONTENU À REMPLIR: Description de la structure du projet frontend]

% Exemple:
% \begin{verbatim}
% src/
%   ├── components/
%   ├── pages/
%   ├── services/
%   ├── styles/
%   └── App.tsx
% \end{verbatim}

\subsection{Composants développés}

% TODO: Décrire les composants principaux
% Points à inclure:
% - Composants réutilisables
% - Composants de pages
% - Gestion de l'état
% - Communication avec l'API

[CONTENU À REMPLIR: Description des composants développés]

\subsection{Captures d'écran des interfaces}

% TODO: Ajouter des captures d'écran
% Points à inclure:
% - Page d'accueil
% - Pages principales
% - Formulaires
% - Tableaux et listes

[CONTENU À REMPLIR: Captures d'écran des interfaces (à remplacer par des figures réelles)]

% Exemple:
% \begin{figure}[H]
%     \centering
%     \includegraphics[width=0.8\textwidth]{images/interface_1.png}
%     \caption{Interface principale de l'application}
%     \label{fig:interface_1}
% \end{figure}

\section{Implémentation du backend}

% TODO: Décrire l'implémentation du backend
% Points à inclure:
% - Structure de l'API
% - Endpoints développés
% - Logique métier
% - Gestion des erreurs

\subsection{Architecture de l'API}

[CONTENU À REMPLIR: Description de l'architecture de l'API]

% Exemple:
% \begin{itemize}
%     \item RESTful API
%     \item Endpoints principaux
%     \item Authentification et autorisations
% \end{itemize}

\subsection{Implémentation des fonctionnalités clés}

% TODO: Détailler l'implémentation des fonctionnalités principales
% Points à inclure:
% - Modules principaux
% - Algorithmes utilisés
% - Logique métier

[CONTENU À REMPLIR: Implémentation des fonctionnalités clés]

% Exemple:
% \begin{verbatim}
% // Exemple de code
% function createUser(userData) {
%     // Implémentation
% }
% \end{verbatim}

\subsection{Gestion de la base de données}

[CONTENU À REMPLIR: Description de la gestion de la base de données]

% Exemple:
% \begin{itemize}
%     \item ORM utilisé: [ORM]
%     \item Migrations: [Outils de migration]
%     \item Seed data: [Description]
% \end{itemize}

\section{Tests et validation}

% TODO: Décrire les stratégies de test
% Points à inclure:
% - Tests unitaires
% - Tests d'intégration
% - Tests de performance
% - Couverture de tests

\subsection{Stratégie de test}

[CONTENU À REMPLIR: Description de la stratégie de test]

% Exemple:
% \begin{itemize}
%     \item Tests unitaires: [Framework]
%     \item Tests d'intégration: [Approche]
%     \item Tests de performance: [Outils]
%     \item Tests d'acceptation: [Méthode]
% \end{itemize}

\subsection{Résultats des tests}

[CONTENU À REMPLIR: Résultats et rapport des tests]

% Exemple:
% \begin{table}[H]
%     \centering
%     \begin{tabular}{|l|c|c|c|}
%         \hline
%         \textbf{Type de test} & \textbf{Nombre} & \textbf{Réussi} & \textbf{Échoué} \\
%         \hline
%         Tests unitaires & 50 & 50 & 0 \\
%         Tests intégration & 20 & 20 & 0 \\
%         Tests performance & 10 & 10 & 0 \\
%         \hline
%     \end{tabular}
%     \caption{Résultats des tests}
%     \label{tab:test_results}
% \end{table}

\section{Déploiement et mise en production}

% TODO: Décrire le processus de déploiement
% Points à inclure:
% - Environnements (dev, staging, prod)
% - Processus de déploiement
% - Gestion des configurations
% - Monitoring et logs

\subsection{Infrastructure de déploiement}

[CONTENU À REMPLIR: Description de l'infrastructure de déploiement]

% Exemple:
% \begin{itemize}
%     \item Serveur de production: [Type]
%     \item Système d'exploitation: [OS]
%     \item Conteneurisation: [Docker/Kubernetes]
%     \item Reverse proxy: [Serveur]
% \end{itemize}

\subsection{Pipeline CI/CD}

[CONTENU À REMPLIR: Description du pipeline CI/CD]

\section{Difficultés rencontrées et solutions}

% TODO: Documenter les difficultés et les solutions
% Points à inclure:
% - Problèmes techniques
% - Défis de conception
% - Solutions apportées
% - Leçons apprises

[CONTENU À REMPLIR: Difficultés rencontrées et solutions apportées]

% Exemple:
% \begin{itemize}
%     \item \textbf{Problème 1:} [Description] - \textbf{Solution:} [Solution]
%     \item \textbf{Problème 2:} [Description] - \textbf{Solution:} [Solution]
%     \item \textbf{Problème 3:} [Description] - \textbf{Solution:} [Solution]
% \end{itemize}

\section*{Conclusion}
\addcontentsline{toc}{section}{Conclusion}

% TODO: Conclure le chapitre 4
% Points à inclure:
% - Synthèse de la réalisation
% - Résultats obtenus
% - État de complétion
% - Transition vers la conclusion générale

[CONTENU À REMPLIR: Conclusion du chapitre 4]
